\section{Final Design}
\label{sec:final-design}

\subsection{Mechanical Architecture}
The selected chassis is a bolted 20\,mm aluminum-extrusion frame carrying four
brushless-hub-motor wheels in a differential-steer arrangement, an IP54
composite payload enclosure with a heated/cooled liner (R-02), and a removable
battery sled accessible from the rear panel without tools (R-12).

\begin{figure}[htbp]
\centering
\begin{tikzpicture}[scale=0.95,every node/.style={font=\small}]
  % chassis outline
  \draw[nexanavy,line width=1.2pt,rounded corners=6pt,fill=fogray]
    (-2.6,-1.4) rectangle (2.6,1.4);
  % payload box
  \draw[steelblue,line width=1pt,rounded corners=3pt,fill=white]
    (-1.5,0.1) rectangle (1.5,1.15);
  \node[color=nexanavy] at (0,0.62) {Insulated Payload Bay};
  % battery sled
  \draw[signalorange,line width=1pt,rounded corners=2pt,fill=white]
    (-2.2,-1.15) rectangle (-0.9,-0.35);
  \node[color=inkgray] at (-1.55,-0.75) {\scriptsize Battery Sled};
  % compute module
  \draw[steelblue,line width=1pt,rounded corners=2pt,fill=white]
    (0.9,-1.15) rectangle (2.2,-0.35);
  \node[color=inkgray] at (1.55,-0.75) {\scriptsize Compute / Sensors};
  % wheels
  \foreach \x/\y in {-2.2/-1.4,2.2/-1.4,-2.2/1.4,2.2/1.4}{
    \draw[nexanavy,fill=inkgray!70] (\x,\y) circle (0.28);
  }
  % LIDAR mast
  \draw[nexanavy,line width=1pt] (0,1.4) -- (0,1.9);
  \draw[signalorange,fill=signalorange] (0,1.9) circle (0.12);
  \node[color=inkgray,anchor=west] at (0.18,1.9) {\scriptsize 2D LIDAR mast};
\end{tikzpicture}
\caption{Top-view schematic of the Nimbus chassis layout (not to scale).}
\label{fig:chassis}
\end{figure}

\subsection{Perception and Navigation Stack}
A 2D LIDAR mounted on a raised mast (Figure~\ref{fig:chassis}) provides the
primary 4\,m-range obstacle field (R-05), fused with wheel odometry and a
forward RGB camera in an extended Kalman filter following the factored
localization approach of \citet{montemerlo2002fastslam}. Path planning runs
a pre-mapped waypoint graph of the three pilot routes with a local
dynamic-window planner \citep{fox1997dynamic} for pedestrian avoidance,
chosen over a vector-field-histogram planner \citep{ulrich2000vfh} for its
lower compute cost on the onboard module; the planner commands a hard stop
whenever an obstacle enters the 4\,m detection envelope while closing speed
exceeds the R-06 stopping-distance budget. Differential-drive kinematics for
the skid-steer controller follow the standard formulation in
\citet{craig2005introduction}.

\subsection{Power and Compute}
A 48\,V, 20\,Ah lithium-iron-phosphate pack (chosen for its stable current
draw at low ambient temperature, unlike the NMC cell it replaced in the fall
prototype) drives four 150\,W hub motors through independent motor
controllers, and feeds a 12\,V rail for the compute module, LIDAR, and
thermal liner via a buck converter bank. See Appendix~\ref{app:bom} for the
full parts and cost accounting behind this subsystem.

\subsection{Design Alternatives Matrix Summary}
Table~\ref{tab:decisions} summarizes the final embodiment decision at each
subsystem boundary, cross-referenced to the trade study in
Section~\ref{sec:alternatives} and the requirement it primarily satisfies.

\begin{table}[htbp]
\centering
\caption{Subsystem design decisions and requirement traceability.}
\label{tab:decisions}
\small
\begin{tabularx}{\textwidth}{@{}lXl@{}}
\toprule
\textbf{Subsystem} & \textbf{Final Decision} & \textbf{Req.} \\
\midrule
Locomotion   & Four-wheel differential skid-steer, low-durometer skid pads & R-04, R-11 \\
Battery      & 48\,V 20\,Ah LiFePO$_4$, tool-free rear sled & R-03, R-12 \\
Perception   & 2D LIDAR (4\,m) + RGB camera, EKF fusion & R-05, R-07 \\
Enclosure    & IP54 composite shell, heated/cooled liner & R-02, R-09 \\
Braking      & Regenerative + electromechanical disc, redundant & R-06 \\
\bottomrule
\end{tabularx}
\end{table}
